\section{Ballistic Phonon NEGF}
\label{sec:res_ballistic}
% ======================================================================

As a heterogeneous test, we model a ballistic Si/Ge interface in the
mass-mismatch approximation, as done in \cite{latour_ab_2017, tian_enhancing_2012, guo_quantum_2020}: bulk-Si force constants in the whole device, with
the Si/Ge distinction entering only through the atomic masses
($m_\mathrm{Si}=28.09$, $m_\mathrm{Ge}=72.64$~amu). At $\Gamma$, the Si
optical maximum is $15.55$~THz (device force-constant set; cf.\ the
$15.12$--$15.37$~THz primitive-cell cross-check of \cref{sec:res_fc}) and the Ge mass-clone gives $9.67$~THz,
consistent with $f_\mathrm{Ge}=f_\mathrm{Si}\sqrt{m_\mathrm{Si}/m_\mathrm{Ge}}$.
On a $20\times20$ transverse mesh, the interfacial conductance at
$300$~K is $G=196$~MW\,m$^{-2}$K$^{-1}$, within $\sim7\%$ of the
literature values ($210$, Latour 2017; $209$, Tian 2012). The
remaining gap is consistent with the coarser transverse mesh and DFT calculations.

% ======================================================================
\section{Force Constants}
\label{sec:res_fc}
% ======================================================================

\paragraph{Method cross-check.} Bulk Silicon harmonic force constants
computed by density-functional perturbation theory (QE \texttt{ph.x}
$\to$ \texttt{q2r.x}) and by finite displacement (phono3py + symfc) on
the same primitive cell agree to $1.6\%$ on the zone-centre optical
frequency ($15.12$ vs $15.37$~THz; \cref{fig:res_fcmethod}). DFPT third-order
% REVIEW(verified): "can hence not be calculated to the same accuracy" was an unsupported accuracy claim; the archived fact is only that d3q rejects the ultrasoft pseudopotential of the harmonic runs ("US not implemented"), so a like-for-like FC3 comparison was not carried out.
constants via \texttt{d3q} require a norm-conserving pseudopotential -- unlike
the ultrasoft one of the harmonic DFPT runs -- so a like-for-like DFPT
cross-check of the third-order constants was not carried out.

\begin{figure}[htb]
  \centering
  \includegraphics[width=0.6\linewidth]{transport_sweeps/fc_method_dispersion_si.pdf}
  \caption{Bulk-silicon phonon dispersion from DFPT versus
  finite-displacement force constants.}
  \label{fig:res_fcmethod}
\end{figure}

\paragraph{Bulk Si $\kappa$.} As an independent check of the force-constant
pipeline and the physics, separate from the NEGF solver, the finite-displacement
FC3 is fed to a phono3py relaxation-time-approximation Boltzmann calculation. The
result is essentially mesh-converged, bulk-silicon
$\kappa(300\,\mathrm{K})=110.7$~W\,m$^{-1}$K$^{-1}$ on an $11^3$ $q$-mesh and
$108.7$ on $19^3$, $\sim27.5\%$ below the experimental $\sim150$, consistent with the
coarse $2\times2\times2$ FC3 supercell (cutoff $5.5$~\AA). Sweeping temperature with
natural-isotope scattering gives a direct comparison to the classic
Glassbrenner--Slack single-crystal
data~\cite{glassbrenner_thermal_1964} (\cref{tab:res_kappaT}): the calculation
reproduces the experimental temperature dependence -- the high-temperature
$\kappa\propto1/T$ fall-off is captured (the conductivity halves between $300$ and
$600$~K) -- and lies
$\sim25$--$30\%$ low near $300$~K, falling to $\sim16$--$18\%$ by $500$--$600$~K, the deficit growing toward
low temperature where long-range cubic and isotope terms matter most. Both effects
are consistent with an under-converged $2\times2\times2$ FC3 supercell.
% REVIEW(verified): the 5x5x5 repeat exists (archive F29: 116.4 W/mK at 19^3 with isotopes); replaced the "is expected to reduce" speculation with the measured result, which closes only ~1/5 of the gap.
A repeat with the better ($5\times5\times5$ VASP) Si FC3 of \cref{sec:res_qfilm}, however, raises
the $19^3$ isotope-scattering value to $\approx116$~W\,m$^{-1}$K$^{-1}$ -- closing only
about a fifth of the gap to experiment, so the FC3 supercell size is not the whole
explanation of the residual deficit.

\begin{table}[t]
  \centering
  \caption{Bulk-Si lattice thermal conductivity vs temperature ($19^3$ mesh, RTA,
    natural isotopes) compared to experimental single-crystal silicon
    (Glassbrenner--Slack).}
  \label{tab:res_kappaT}
  \begin{tabular}{rrrr}
    \toprule
    $T$ (K) & $\kappa_\text{calc}$ (W/mK) & $\kappa_\text{exp}$ (W/mK) & ratio\\
    \midrule
    $200$ & $179.6$ & ${\sim}266$ & $0.68$\\
    % REVIEW(verified): 108.7/150 = 0.72 (the archived 0.73 traces to a finer experimental value than the tabulated ~150); table made self-consistent.
    $300$ & $108.7$ & ${\sim}150$ & $0.72$\\
    $400$ & $78.7$  & ${\sim}98$  & $0.80$\\
    $500$ & $62.0$  & ${\sim}76$  & $0.82$\\
    $600$ & $51.3$  & ${\sim}61$  & $0.84$\\
    \bottomrule
  \end{tabular}
\end{table}

\paragraph{FC3 compression.} The tensor-decomposition methods of
\cref{sub:fc3_compression} are applied to the device FC3 of the
\texttt{d11a} nanowire ($39.4\times10^6$ dense parameters) and the
compressed tensors fed back through the full anharmonic transport solve
(\cref{fig:res_decomp}). We observe that, first, the transport
observable converges faster than the Frobenius residual the fitters
minimise: the symmetric ans\"atze (INDSCAL, Waring) reach a $\sim15\%$
conductance error at rank $16$ ($\sim3600$--$4500\times$ fewer
parameters) while their Frobenius residual is still $52$--$67\%$. The
relevant fitting target for phonon NEGF is therefore a transport
residual, not the tensor norm. Second, the current compression approach
has distinct failure modes: the
matricised SVD is both far denser (its factor count is up to ${\sim}110\times$
the symmetric methods' at equal rank) and unphysical at higher rank, with its
rank-8 fit violating heat-flow conservation and its rank-16 fit giving a
negative $G_\mathrm{anh}$. Conversely, the Tucker (HOSVD) and INDSCAL fits at
the lowest ranks collapse to the bare ballistic value. Only the symmetric INDSCAL/Waring families degrade
gracefully.

\begin{figure}[htb]
  \centering
  \includegraphics[width=0.48\linewidth]{transport_sweeps/d11a_decomp_ganh.pdf}\hfill
  \includegraphics[width=0.48\linewidth]{transport_sweeps/d11a_decomp_conservation.pdf}
  \caption{FC3 tensor-decomposition transport quality on the \texttt{d11a}
  nanowire: anharmonic conductance (left) and heat-flow conservation
  residual (right) versus FC3 parameter count, for five decomposition
  families against the dense reference.}
  \label{fig:res_decomp}
\end{figure}

% ======================================================================
